\section{Introduction}

A bus can spend the same minute in three different ways: serving customers sooner, holding to regulate spacing, or returning to cover the next scheduled trip. Allocating more trips to a peak helps only if vehicles can actually depart. Likewise, a hold that improves headways at one stop can make the bus unavailable when the opposite terminal needs it. We ask whether joint dispatch and holding improve passenger service under a fixed physical fleet and trip roster, and how much holding adds once a learned dispatch policy is already in place.

Gap-feedback holding can moderate bunching but incurs time for riders already aboard \citep{daganzo2009headway,xuan2011dynamic,berrebi2018comparing}. Optimization-based methods account for capacity, waiting and holding costs \citep{eberlein2001holding,hickman2001analytic,delgado2012much,sanchez2016realtime}. Dispatch research has studied how to shift a fixed service supply toward higher demand \citep{newell1971dispatching}. \citet{sanchez2017schedulefree} examine real-time headway-managed dispatch and stop-skipping under operator resource constraints. Most directly, \citet{gkiotsalitis2021rescheduling} jointly optimizes departure rescheduling and bus holding with explicit original-timetable constraints, evaluating excess waiting against in-vehicle time in an operator-data simulation. Our emphasis is different: in a closed physical fleet, we measure what the same dispatch policy gains or loses when changes to holding alter the pool of buses available for subsequent departures. A same-authority rolling comparator extends the comparison beyond fixed timetables.

This paper contributes (i) an executable terminal-to-stop coordination rule that caps discretionary holding with a trip budget, (ii) paired comparisons that separate the effect of the complete learned controller, the choice of actor inputs, and the use of stop holding under a common upper policy, and (iii) a vehicle-level account of a circulation failure. The point is not to show that frequency separation guarantees better operations. It is to establish which observed gains survive a realistic base comparison and which disappear when a component is changed.
